\begin{exercises}
  \item 求解该方程组。
    再求解相应的齐次方程组。
    \begin{equation*}
      \begin{linsys}{4}
        x  &+ &y  &- &2z &= &0 \\
        x  &- &y  &  &   &= &-3 \\
        3x &- &y  &- &2z &= &-6  \\
           &  &2y &- &2z &= &3
      \end{linsys}
   \end{equation*}
   \begin{answer}
     下面的化简求解了该方程组。
     \begin{align*}
        \begin{amat}{3}
          1 &1   &-2 &0 \\
          1 &-1  &0  &3  \\
          3 &-1  &-2 &-6 \\
          0 &2   &-2 &3
        \end{amat}
        &\grstep[-3\rho_1+\rho_3]{-\rho_1+\rho_2}
        \begin{amat}{3}
          1 &1   &-2 &0 \\
          0 &-2  &2  &-3  \\
          0 &-4  &4  &-6 \\
          0 &2   &-2 &3
        \end{amat}                                 \\
        &\grstep[\rho_2+\rho_4]{-2\rho_2+\rho_3}
        \begin{amat}{3}
          1 &1   &-2 &0 \\
          0 &-2  &2  &-3  \\
          0 &0   &0  &0 \\
          0 &0   &0  &0
        \end{amat}
      \end{align*}
      解集如下。
      \begin{equation*}
        \set{\colvec{-3/2 \\ 3/2 \\ 0}
             +\colvec{1 \\ 1 \\ 1}z
             \suchthat z\in\Re}
      \end{equation*}
      类似地，我们可以化简相应的齐次方程组
      \begin{align*}
        \begin{amat}{3}
          1 &1   &-2 &0 \\
          1 &-1  &0  &0  \\
          3 &-1  &-2 &0 \\
          0 &2   &-2 &0
        \end{amat}
        &\grstep[-3\rho_1+\rho_3]{-\rho_1+\rho_2}
        \begin{amat}{3}
          1 &1   &-2 &0 \\
          0 &-2  &2  &0  \\
          0 &-4  &4  &0 \\
          0 &2   &-2 &0
        \end{amat}                                         \\
        &\grstep[\rho_2+\rho_4]{-2\rho_2+\rho_3}
        \begin{amat}{3}
          1 &1   &-2 &0 \\
          0 &-2  &2  &0  \\
          0 &0   &0  &0 \\
          0 &0   &0  &0
        \end{amat}
      \end{align*}
      得到它的解集。
      \begin{equation*}
        \set{
             \colvec{1 \\ 1 \\ 1}z
             \suchthat z\in\Re}
      \end{equation*}
    \end{answer}
  \recommended \item
    求解每个方程组。
    用向量表示解集。
    指出一个特解以及齐次方程组的解集。
    % (These systems also appear in
    % Exercise~2.\nearbyexercise{exer:SolveInMatrixNotation}.)
    \begin{exparts*}
      \partsitem \( \begin{linsys}[t]{2}
                  3x  &+  &6y  &=  &18  \\
                   x  &+  &2y  &=  &6
                    \end{linsys}  \)
      \partsitem \( \begin{linsys}[t]{2}
                   x  &+  &y   &=  &1  \\
                   x  &-  &y   &=  &-1
                    \end{linsys}  \)
      \partsitem \( \begin{linsys}[t]{3}
                   x_1  &   &     &+  &x_3   &=  &4  \\
                   x_1  &-  &x_2  &+  &2x_3  &=  &5  \\
                  4x_1  &-  &x_2  &+  &5x_3  &=  &17
                    \end{linsys}  \)
      \partsitem \( \begin{linsys}[t]{3}
                   2a   &+  &b    &-  &c     &=  &2  \\
                   2a   &   &     &+  &c     &=  &3  \\
                    a   &-  &b    &   &      &=  &0
                    \end{linsys}  \)
      \partsitem \( \begin{linsys}[t]{4}
                     x  &+  &2y   &-   &z   &    &    &=  &3  \\
                    2x  &+  &y    &    &    &+   &w   &=  &4  \\
                     x  &-  &y    &+   &z   &+   &w   &=  &1
                    \end{linsys}  \)
      \partsitem \( \begin{linsys}[t]{4}
                     x  &   &     &+   &z   &+   &w   &=  &4  \\
                    2x  &+  &y    &    &    &-   &w   &=  &2  \\
                    3x  &+  &y    &+   &z   &    &    &=  &7
                    \end{linsys}  \)
    \end{exparts*}
    \begin{answer}
      这几题的具体计算参见前一小节的答案。

      \begin{exparts}
        \partsitem
          解集如下。
          \begin{equation*}
            S=\set{\colvec[r]{6 \\ 0}+\colvec[r]{-2 \\ 1}y
              \suchthat y\in\Re}
          \end{equation*}
          下面是特解以及相应的齐次方程组的解集。
          \begin{equation*}
            \colvec[r]{6 \\ 0}
              \quad\text{与}\quad
            \set{\colvec[r]{-2 \\ 1}y
              \suchthat y\in\Re}
          \end{equation*}
          \textit{注。}
          这里可能有两点混淆。
          首先，上面给出的集合 $S$ 等于这个集合
          \begin{equation*}
            T=\set{\colvec[r]{4 \\ 1}+\colvec[r]{-2 \\ 1}y
              \suchthat y\in\Re}
          \end{equation*}
          因为这两个集合含有相同的元素。
          对于“什么是特解？”这一问题，下面这些都是正确的答案：
          \begin{equation*}
            \colvec{6 \\ 0},\quad
            \colvec{4 \\ 1},\quad
            \colvec{2 \\ 2},\quad
            \colvec{1 \\ 2.5}
          \end{equation*}
          第二点混淆是，我们在集合中使用的字母无关紧要。
          这个集合也等于 $S$。
          \begin{equation*}
            U=\set{\colvec[r]{6 \\ 0}+\colvec[r]{-2 \\ 1}u
                   \suchthat u\in\Re}
          \end{equation*}
        \partsitem
          解集如下。
          \begin{equation*}
            \set{\colvec[r]{0 \\ 1} }
          \end{equation*}
          下面是一个特解，
          以及相应的齐次方程组的解集。
          \begin{equation*}
            \colvec[r]{0 \\ 1}
              \qquad
            \set{\colvec[r]{0 \\ 0} }
          \end{equation*}
        \partsitem
          解集是无穷的。
          \begin{equation*}
            \set{\colvec[r]{4 \\ -1 \\ 0}+\colvec[r]{-1 \\ 1 \\ 1}x_3
              \suchthat x_3\in\Re}
          \end{equation*}
          下面是一个特解以及相应的齐次方程组的解集。
          \begin{equation*}
            \colvec[r]{4 \\ -1 \\ 0}
              \qquad
            \set{\colvec[r]{-1 \\ 1 \\ 1}x_3
              \suchthat x_3\in\Re}
          \end{equation*}
        \partsitem
          解集是单元素集。
          \begin{equation*}
            \set{\colvec[r]{1 \\ 1 \\ 1}}
          \end{equation*}
          下面是一个特解以及相应的齐次方程组的解集。
          \begin{equation*}
            \colvec[r]{1 \\ 1 \\ 1}
              \qquad
            \set{\colvec[r]{0 \\ 0 \\ 0}}
          \end{equation*}
        \partsitem
          解集是无穷的。
          \begin{equation*}
            \set{\colvec[r]{5/3 \\ 2/3 \\ 0 \\ 0}
                 +\colvec[r]{-1/3 \\ 2/3 \\ 1 \\ 0}z
                 +\colvec[r]{-2/3 \\ 1/3 \\ 0 \\ 1}w
                 \suchthat z,w\in\Re}
          \end{equation*}
          下面是一个特解以及相应的齐次方程组的解集。
          \begin{equation*}
            \colvec[r]{5/3 \\ 2/3 \\ 0 \\ 0}
              \qquad
            \set{\colvec[r]{-1/3 \\ 2/3 \\ 1 \\ 0}z
                 +\colvec[r]{-2/3 \\ 1/3 \\ 0 \\ 1}w
                 \suchthat z,w\in\Re}
          \end{equation*}
        \partsitem 该方程组的解集为空。
          因此，不存在特解。
          相应的齐次方程组的解集如下。
          \begin{equation*}
            \set{\colvec[r]{-1 \\ 2 \\ 1 \\ 0}z
                 +\colvec[r]{-1 \\ 3 \\ 0 \\ 1}w
                 \suchthat z,w\in\Re}
          \end{equation*}
      \end{exparts}
    \end{answer}
  \item
    求解每个方程组，用向量记法给出解集。
    指出一个特解以及齐次方程组的解。
    \begin{exparts*}
      \partsitem \( \begin{linsys}[t]{3}
                  2x  &+  &y  &-  &z  &=  &1  \\
                  4x  &-  &y  &   &   &=  &3
                  \end{linsys}  \)
      \partsitem \( \begin{linsys}[t]{4}
                   x  &   &   &-  &z  &   &   &=  &1  \\
                      &   &y  &+  &2z &-  &w  &=  &3  \\
                   x  &+  &2y &+  &3z &-  &w  &=  &7
                   \end{linsys}  \)
      \partsitem \( \begin{linsys}[t]{4}
                   x  &-  &y  &+  &z  &   &   &=  &0  \\
                      &   &y  &   &   &+  &w  &=  &0  \\
                  3x  &-  &2y &+  &3z &+  &w  &=  &0  \\
                      &   &-y &   &   &-  &w  &=  &0
                  \end{linsys}  \)
      \partsitem \( \begin{linsys}[t]{5}
                   a  &+  &2b &+  &3c &+  &d  &-  &e  &=  &1  \\
                  3a  &-  &b  &+  &c  &+  &d  &+  &e  &=  &3
                  \end{linsys}  \)
    \end{exparts*}
    \begin{answer}
    前一小节的答案给出了行变换。
    这里的每个答案只列出解集、特解和齐次解。
    \begin{exparts}
      \partsitem
        解集如下。
        \begin{equation*}
          \set{\colvec[r]{2/3 \\ -1/3 \\ 0}
               +\colvec[r]{1/6 \\ 2/3 \\ 1}z
              \suchthat z\in\Re}
        \end{equation*}
        下面是一个特解以及相应的齐次方程组的解集。
        \begin{equation*}
          \colvec[r]{2/3 \\ -1/3 \\ 0}
            \qquad
        \set{\colvec[r]{1/6 \\ 2/3 \\ 1}z
            \suchthat z\in\Re}
        \end{equation*}
      \partsitem
        解集是无穷的。
        \begin{equation*}
          \set{\colvec[r]{1 \\ 3 \\ 0 \\ 0}
               +\colvec[r]{1 \\-2 \\ 1 \\ 0}z
              \suchthat z\in\Re}
        \end{equation*}
        下面是
        一个特解以及相应的齐次方程组的解集。
        \begin{equation*}
          \colvec[r]{1 \\ 3 \\ 0 \\ 0}
            \qquad
          \set{\colvec[r]{1 \\ -2 \\ 1 \\ 0}z
              \suchthat z\in\Re}
        \end{equation*}
      \partsitem
        解集如下。
        \begin{equation*}
          \set{\colvec[r]{0 \\ 0 \\ 0 \\ 0}
               +\colvec[r]{-1 \\ 0 \\ 1 \\ 0}z
               +\colvec[r]{-1 \\ -1 \\ 0 \\ 1}w
              \suchthat z,w\in\Re}
        \end{equation*}
        下面是
        一个特解以及相应的齐次方程组的解集。
        \begin{equation*}
          \colvec[r]{0 \\ 0 \\ 0 \\ 0}
            \qquad
          \set{\colvec[r]{-1 \\ 0 \\ 1 \\ 0}z
               +\colvec[r]{-1 \\ -1 \\ 0 \\ 1}w
              \suchthat z,w\in\Re}
        \end{equation*}
      \partsitem
        解集如下。
        \begin{equation*}
          \set{\colvec[r]{1 \\ 0 \\ 0 \\ 0 \\ 0}
               +\colvec[r]{-5/7 \\ -8/7 \\ 1 \\ 0 \\ 0}c
               +\colvec[r]{-3/7 \\ -2/7 \\ 0 \\ 1 \\ 0}d
               +\colvec[r]{-1/7 \\ 4/7 \\ 0 \\ 0 \\ 1}e
              \suchthat c,d,e\in\Re}
        \end{equation*}
        下面是
        一个特解以及相应的齐次方程组的解集。
        \begin{equation*}
          \colvec[r]{1 \\ 0 \\ 0 \\ 0 \\ 0}
            \qquad
          \set{\colvec[r]{-5/7 \\ -8/7 \\ 1 \\ 0 \\ 0}c
               +\colvec[r]{-3/7 \\ -2/7 \\ 0 \\ 1 \\ 0}d
               +\colvec[r]{-1/7 \\ 4/7 \\ 0 \\ 0 \\ 1}e
              \suchthat c,d,e\in\Re}
        \end{equation*}
    \end{exparts}
   \end{answer}
  \recommended \item
    对于方程组
    \begin{equation*}
      \begin{linsys}{4}
       2x  &-  &y  &   &    &-  &w  &=  &3  \\
           &   &y  &+  &z   &+  &2w &=  &2  \\
        x  &-  &2y &-  &z   &   &   &=  &-1
      \end{linsys}
    \end{equation*}
    下列向量中哪些可以用作某个通解中的特解部分？
    \begin{exparts*}
      \partsitem   \( \colvec[r]{0 \\ -3 \\ 5 \\ 0} \)
      \partsitem   \( \colvec[r]{2 \\ 1 \\ 1 \\ 0} \)
      \partsitem   \( \colvec[r]{-1 \\ -4 \\ 8 \\ -1} \)
    \end{exparts*}
    \begin{answer}
      直接把它们代入，看它们是否满足全部三个方程。
      \begin{exparts}
        \partsitem 否。
        \partsitem 是。
        \partsitem 是。
      \end{exparts}
    \end{answer}
  \recommended \item
    \nearbylemma{th:GenEqPartHomo} 表明 $\vec{p}$ 可以取任意一个特解。
    若可能，求出这个方程组的通解
    \begin{equation*}
      \begin{linsys}{4}
        x  &-  &y  &   &    &+  &w  &=  &4  \\
       2x  &+  &3y &-  &z   &   &   &=  &0  \\
           &   &y  &+  &z   &+  &w  &=  &4
      \end{linsys}
    \end{equation*}
    使之以给定的向量作为其特解。
    \begin{exparts*}
      \partsitem   \( \colvec[r]{0 \\ 0 \\ 0 \\ 4} \)
      \partsitem   \( \colvec[r]{-5 \\ 1 \\ -7 \\ 10} \)
      \partsitem   \( \colvec[r]{2 \\ -1 \\ 1 \\ 1} \)
    \end{exparts*}
    \begin{answer}
      对相应的齐次方程组施行高斯消元法
      \begin{align*}
        \begin{amat}[r]{4}
           1  &-1  &0  &1  &0  \\
           2  &3   &-1 &0  &0  \\
           0  &1   &1  &1  &0
        \end{amat}
        &\grstep{-2\rho_1+\rho_2}
        \begin{amat}[r]{4}
           1  &-1  &0  &1  &0  \\
           0  &5   &-1 &-2 &0  \\
           0  &1   &1  &1  &0
        \end{amat}                                \\
        &\grstep{-(1/5)\rho_2+\rho_3}
        \begin{amat}[r]{4}
           1  &-1  &0  &1  &0  \\
           0  &5   &-1 &-2 &0  \\
           0  &0   &6/5&7/5&0
        \end{amat}
      \end{align*}
      得到这个解是齐次问题的解。
      \begin{equation*}
        \set{\colvec[r]{-5/6 \\ 1/6 \\ -7/6 \\ 1}w\suchthat w\in\Re}
      \end{equation*}
      \begin{exparts}
        \partsitem 该向量确实是一个特解，因此所求的
          通解如下。
          \begin{equation*}
            \set{\colvec[r]{0 \\ 0 \\ 0 \\ 4}+
                 \colvec[r]{-5/6 \\ 1/6 \\ -7/6 \\ 1}w\suchthat w\in\Re}
          \end{equation*}
        \partsitem 该向量是一个特解，因此所求的
          通解如下。
          \begin{equation*}
            \set{\colvec[r]{-5 \\ 1 \\ -7 \\ 10}+
                 \colvec[r]{-5/6 \\ 1/6 \\ -7/6 \\ 1}w\suchthat w\in\Re}
          \end{equation*}
        \partsitem 该向量不是该方程组的解，因为
          它不满足第三个方程。
          这样的通解不存在。
      \end{exparts}
    \end{answer}
  \item
     下面两个中，一个是非奇异的而另一个是奇异的。
     请指出哪个是哪个？
     \begin{exparts*}
       \partsitem $\begin{mat}[r]
           1  &3   \\
           4  &-12
         \end{mat}$
       \partsitem $\begin{mat}[r]
           1  &3  \\
           4  &12
         \end{mat}$
     \end{exparts*}
     \begin{answer}
       第一个非奇异而第二个奇异。
       直接用高斯消元法，看阶梯形结果的
       对角线上每个位置的数是否都非 $0$。
     \end{answer}
  \recommended \item
    奇异还是非奇异？
    \begin{exparts*}
      \partsitem \(
        \begin{mat}[r]
          1  &2  \\
          1  &3
        \end{mat}   \)
      \partsitem \(
        \begin{mat}[r]
          1  &2  \\
         -3  &-6
        \end{mat}   \)
      \partsitem \(
        \begin{mat}[r]
          1  &2  &1  \\
          1  &3  &1
        \end{mat}   \)
      \partsitem \(
        \begin{mat}[r]
          1  &2  &1  \\
          1  &1  &3  \\
          3  &4  &7
        \end{mat}   \)
      \partsitem \(
        \begin{mat}[r]
          2  &2  &1  \\
          1  &0  &5  \\
         -1  &1  &4
        \end{mat}   \)
    \end{exparts*}
    \begin{answer}
      \begin{exparts}
      \partsitem 非奇异：
        \begin{equation*}
          \grstep{-\rho_1+\rho_2}
          \begin{mat}[r]
            1  &2  \\
            0  &1
          \end{mat}
        \end{equation*}
        最终每行都含有一个首主元。
      \partsitem 奇异：
        \begin{equation*}
          \grstep{3\rho_1+\rho_2}
          \begin{mat}[r]
            1  &2  \\
            0  &0
          \end{mat}
        \end{equation*}
        最终第 \( 2 \) 行没有首主元。
      \partsitem 都不是。
        矩阵必须是方阵，这两个词才适用。
      \partsitem 奇异。
      \partsitem 非奇异。
     \end{exparts}
    \end{answer}
  \recommended \item
    给定的向量是否属于由给定的集合生成的集合？
      \begin{exparts}
        \partsitem \( \colvec[r]{2 \\ 3}, \)\
          \( \set{\colvec[r]{1 \\ 4},
                \colvec[r]{1 \\ 5}} \)
        \partsitem \( \colvec[r]{-1 \\ 0 \\ 1}, \)\
          \( \set{\colvec[r]{2 \\ 1 \\ 0},
                \colvec[r]{1 \\ 0 \\ 1}} \)
        \partsitem \( \colvec[r]{1 \\ 3 \\ 0}, \)\
          \( \set{\colvec[r]{1 \\ 0 \\ 4},
                \colvec[r]{2 \\ 1 \\ 5},
                \colvec[r]{3 \\ 3 \\ 0},
                \colvec[r]{4 \\ 2 \\ 1}} \)
        \partsitem \( \colvec[r]{1 \\ 0 \\ 1 \\ 1}, \)\
          \( \set{\colvec[r]{2 \\ 1 \\ 0 \\ 1},
                \colvec[r]{3 \\ 0 \\ 0 \\ 2}} \)
      \end{exparts}
      \begin{answer}
        在每种情形下，我们都要判断该向量是否为集合中那些向量的
        线性组合。
        \begin{exparts}
          \partsitem 是。
            求解
            \begin{equation*}
              c_1\colvec[r]{1 \\ 4}+c_2\colvec[r]{1 \\ 5}=\colvec[r]{2 \\ 3}
            \end{equation*}
            结合
            \begin{equation*}
              \begin{amat}[r]{2}
                1  &1  &2  \\
                4  &5  &3
              \end{amat}
              \grstep{-4\rho_1+\rho_2}
              \begin{amat}[r]{2}
                1  &1  &2  \\
                0  &1  &-5
              \end{amat}
            \end{equation*}
            可以得出存在给出该组合的 $c_1$ 和 $c_2$。
          \partsitem 否。
            化简
            \begin{equation*}
              \begin{amat}[r]{2}
                2  &1  &-1 \\
                1  &0  &0  \\
                0  &1  &1
              \end{amat}
              \grstep{-(1/2)\rho_1+\rho_2}
              \begin{amat}[r]{2}
                2  &1     &-1 \\
                0  &-1/2  &1/2  \\
                0  &1     &1
              \end{amat}
              \grstep{2\rho_2+\rho_3}
              \begin{amat}[r]{2}
                2  &1     &-1 \\
                0  &-1/2  &1/2  \\
                0  &0     &2
              \end{amat}
            \end{equation*}
            表明
            \begin{equation*}
              c_1\colvec[r]{2 \\ 1 \\ 0}+c_2\colvec[r]{1 \\ 0 \\ 1}
                =\colvec[r]{-1 \\ 0 \\ 1}
            \end{equation*}
            无解。
          \partsitem 是。
            化简
            \begin{align*}
              \begin{amat}[r]{4}
                1  &2  &3  &4  &1  \\
                0  &1  &3  &2  &3  \\
                4  &5  &0  &1  &0
              \end{amat}
              &\grstep{-4\rho_1+\rho_3}
              \begin{amat}[r]{4}
                1  &2  &3  &4  &1  \\
                0  &1  &3  &2  &3  \\
                0  &-3 &-12&-15&-4
              \end{amat}                   \\
              &\grstep{3\rho_2+\rho_3}
              \begin{amat}[r]{4}
                1  &2  &3  &4  &1  \\
                0  &1  &3  &2  &3  \\
                0  &0  &-3 &-9 &5
              \end{amat}
            \end{align*}
            表明有无穷多种方式
            \begin{equation*}
              \set{\colvec[r]{c_1 \\ c_2 \\ c_3 \\ c_4}=
                   \colvec[r]{-10 \\ 8 \\ -5/3 \\ 0}+
                   \colvec[r]{-9 \\ 7 \\ -3 \\ 1}c_4
                    \suchthat c_4\in\Re}
            \end{equation*}
            来写出这个组合。
            \begin{equation*}
              \colvec[r]{1 \\ 3 \\ 0}=
              c_1\colvec[r]{1 \\ 0 \\ 4}+
              c_2\colvec[r]{2 \\ 1 \\ 5}+
              c_3\colvec[r]{3 \\ 3 \\ 0}+
              c_4\colvec[r]{4 \\ 2 \\ 1}
            \end{equation*}
          \partsitem 否。
            看第三个分量。
        \end{exparts}
      \end{answer}
  \item
     证明：系数矩阵非奇异的任何线性方程组都有解，并且解是唯一的。
    \begin{answer}
        由于系数矩阵是非奇异的，高斯消元法
        结束时得到的阶梯形中每个变量都在某个方程中居于首位。
        回代给出唯一解。

      （另一种看出解唯一的方法是注意：
      当系数矩阵非奇异时，由定义，
      相应的齐次方程组有唯一解。
      由于通解是一个特解与
      每个齐次解的和，所以通解
      至多只有一个元素。）
     \end{answer}
  \item
    在引理
    \nearbylemma{le:HomoSltnSpanVecs}
    的证明中，如果没有非 \( 0=0 \) 的方程，会发生什么？
    \begin{answer}
      此时解集是整个 \( \Re^n \)，而且我们可以
      用所要求的形式表示它。
      \begin{equation*}
        \set{c_1\colvec[r]{1 \\ 0 \\ \vdotswithin{1} \\ 0}
             +c_2\colvec[r]{0 \\ 1 \\ \vdotswithin{0} \\ 0}
             +\cdots
             +c_n\colvec[r]{0 \\ 0 \\ \vdotswithin{0} \\ 1}
             \suchthat c_1,\ldots,c_n\in\Re}
      \end{equation*}
     \end{answer}
  \recommended \item
    证明：若 \( \vec{s} \) 和 \( \vec{t} \)
    满足某个齐次方程组，则下面这些向量也满足它。
    \begin{exparts*}
      \partsitem \( \vec{s}+\vec{t} \)
      \partsitem \( 3\vec{s} \)
      \partsitem \( k\vec{s}+m\vec{t} \)，其中 \( k,m\in\Re \)
    \end{exparts*}
    这个论证有什么问题：“这三个结果表明，如果齐次
    方程组有一个解，那么它就有许多个解\Dash 解的
    任意倍数是另一个解，任意一些解的和也是
    一个解\Dash 因此不存在恰好有一个解的
    齐次方程组。”？
    \begin{answer}
      设 \( \vec{s},\vec{t}\in\Re^n \)，并把它们写成如下形式。
      \begin{equation*}
        \vec{s}=\colvec{s_1 \\ \vdotswithin{s_1} \\ s_n}
          \qquad
        \vec{t}=\colvec{t_1 \\ \vdotswithin{t_1} \\ t_n}
      \end{equation*}
      再设 \( a_{i,1}x_1+\cdots+a_{i,n}x_n=0 \) 是齐次方程组中的
      第 \( i \) 个方程。
      \begin{exparts}
        \partsitem 验证很容易。
          \begin{multline*}
            a_{i,1}(s_1+t_1)+\cdots+a_{i,n}(s_n+t_n)      \\
            =
            (a_{i,1}s_1+\cdots+a_{i,n}s_n)
            +(a_{i,1}t_1+\cdots+a_{i,n}t_n)
            =
            0+0
          \end{multline*}
        \partsitem 这与前一个类似。
          \begin{equation*}
            a_{i,1}(3s_1)+\cdots+a_{i,n}(3s_n)
            =3(a_{i,1}s_1+\cdots+a_{i,n}s_n)
            =3\cdot 0=0
          \end{equation*}
        \partsitem 这个也不难多少。
          \begin{multline*}
            a_{i,1}(ks_1+mt_1)+\cdots+a_{i,n}(ks_n+mt_n)  \\
            =
            k(a_{i,1}s_1+\cdots+a_{i,n}s_n)
            +m(a_{i,1}t_1+\cdots+a_{i,n}t_n)
            =
            k\cdot 0+m\cdot 0
          \end{multline*}
      \end{exparts}
     这个论证的错误在于，只涉及
     零向量的任何线性组合都会得到零向量。
   \end{answer}
  \item
    证明：如果一个系数和常数全为有理数的方程组
    有解，那么它至少有一个全为有理数的解。
    它必定有无穷多个解吗？
    \begin{answer}
      先看证明。

      若方程组的系数全为有理数，则
      带回代的高斯消元法将只用到有理数（例如，
      \( -(m/n)\rho_i+\rho_j \)）。
      于是，我们可以只用有理数
      作为每个向量的分量来表示解集。
      因此，特解全为有理数。

      至于无穷多解的问题，
      有理向量解有无穷多个，当且仅当
      相应的齐次方程组有无穷多个
      实向量解。
      这是因为把任意参数取成有理数都会产生一个
      全为有理数的解。
   \end{answer}
\end{exercises}

\endinput
